\documentclass[a4paper,12pt,twoside,leqno]{article}
\usepackage[marginratio={5:5, 5:5}, textwidth=150mm, ]{geometry}


\pagenumbering{gobble}
\usepackage{amssymb}
\usepackage{xcolor}
\usepackage{amsmath}
\usepackage{hyperref}
\usepackage{tikz}
\usetikzlibrary{quotes}
\usetikzlibrary{automata, arrows.meta, positioning}
\usepackage{tikz-cd}
\title{Notes on "Arrows, Structures and Functors: Categorical Imperative"}
\author{}
\date{}
\usepackage[most]{tcolorbox}

\newtcolorbox[auto counter, number within=subsection]{theorem}[2][]{%
  colback=blue!5!white,
  colframe=red!25!black,
  fonttitle=\bfseries,
  title=Proposition~\thetcbcounter,
}


\begin{document}
\maketitle
\section{Learning to think with arrows}
\subsection{Epimorphisms and Monomorphisms}

- A function $f: A \rightarrow B$, with \textbf{domain} $A$ and \textbf{codomain} $B$, is \textbf{onto} or \textbf{surjective} if every element $b$ (or $f(a)$) in the codomain $B$ is mapped by at least one element $a$ in the domain $A$. So for instance, $f: \mathbb{Z} \rightarrow \mathbb{Z}: n \mapsto n + 1$ is onto. Another example of such is the identity mapping, e.g. $id_{A}: A \rightarrow A$. An example of a function which isn't onto is for instance $f: \mathbb{Z} \rightarrow \mathbb{Z}: n \mapsto n^2$, since all negative integers in the codomain aren't mapped by any element in the domain.\\
- Two functions are said to be the same if they have the same domain and codomain, and yield the same value given a same argument.

\begin{theorem}

If a map $f: A \rightarrow B$ is onto, then whenever two maps $g, h: B \rightarrow C$ satisfy $g \circ f = h \circ f$, $g$ and $h$ must be equal.

\end{theorem}

\begin{theorem}

\textbf{DEF.} A function $f$ is an \textbf{epimorphism} if $g \circ f = h \circ f$ always implies $g = h$.

\end{theorem}

\begin{theorem}

A function $f: A \rightarrow B$ is onto iff it is an epimorphism.

\end{theorem}
\end{document}
